\chapter{Appendici}

\section{Elenco completo delle feature del dataset IBM HR}
\label{app:features}

La Tabella~\ref{tab:features_complete} riporta le 35 feature originali del dataset IBM HR Employee Attrition, con indicazione del tipo, del range o dei valori ammessi e di una breve descrizione. Le feature contrassegnate con $\dagger$ sono state rimosse in fase di preprocessing in quanto costanti o identificative.

\begin{longtable}{p{4.2cm} c p{3.2cm} p{4.5cm}}
    \caption{Feature del dataset IBM HR Employee Attrition.}
    \label{tab:features_complete} \\
    \toprule
    Feature & Tipo & Range/Valori & Descrizione \\
    \midrule
    \endfirsthead
    \toprule
    Feature & Tipo & Range/Valori & Descrizione \\
    \midrule
    \endhead
    \midrule
    \multicolumn{4}{r}{\textit{Continua nella pagina successiva}} \\
    \bottomrule
    \endfoot
    \bottomrule
    \endlastfoot
    Age & int & [18, 60] & Età del dipendente \\
    Attrition & cat & Yes, No & Variabile target: abbandono \\
    BusinessTravel & cat & Non-Travel, Rarely, Frequently & Frequenza dei viaggi di lavoro \\
    DailyRate & int & [102, 1499] & Tariffa giornaliera \\
    Department & cat & HR, R\&D, Sales & Dipartimento di appartenenza \\
    DistanceFromHome & int & [1, 29] & Distanza casa-lavoro (km) \\
    Education & int & [1, 5] & Livello di istruzione (1=sotto diploma, 5=dottorato) \\
    EducationField & cat & 6 valori & Campo di studi \\
    EmployeeCount$^\dagger$ & int & 1 & Costante (rimossa) \\
    EmployeeNumber$^\dagger$ & int & [1, 2068] & Identificativo (rimosso) \\
    EnvironmentSatisfaction & int & [1, 4] & Soddisfazione ambiente (1=bassa, 4=alta) \\
    Gender & cat & Female, Male & Genere \\
    HourlyRate & int & [30, 100] & Tariffa oraria \\
    JobInvolvement & int & [1, 4] & Coinvolgimento nel lavoro \\
    JobLevel & int & [1, 5] & Livello gerarchico \\
    JobRole & cat & 9 valori & Ruolo ricoperto \\
    JobSatisfaction & int & [1, 4] & Soddisfazione lavorativa \\
    MaritalStatus & cat & Divorced, Married, Single & Stato civile \\
    MonthlyIncome & int & [1009, 19999] & Reddito mensile (\$) \\
    MonthlyRate & int & [2094, 26999] & Tariffa mensile \\
    NumCompaniesWorked & int & [0, 9] & Numero di aziende precedenti \\
    Over18$^\dagger$ & cat & Y & Costante (rimossa) \\
    OverTime & cat & Yes, No & Presenza di straordinari \\
    PercentSalaryHike & int & [11, 25] & Incremento salariale (\%) \\
    PerformanceRating & int & [3, 4] & Valutazione della prestazione \\
    RelationshipSatisfaction & int & [1, 4] & Soddisfazione relazionale \\
    StandardHours$^\dagger$ & int & 80 & Costante (rimossa) \\
    StockOptionLevel & int & [0, 3] & Livello di stock option \\
    TotalWorkingYears & int & [0, 40] & Anni totali di esperienza \\
    TrainingTimesLastYear & int & [0, 6] & Sessioni di formazione nell'ultimo anno \\
    WorkLifeBalance & int & [1, 4] & Equilibrio vita-lavoro \\
    YearsAtCompany & int & [0, 40] & Anni nell'azienda attuale \\
    YearsInCurrentRole & int & [0, 18] & Anni nel ruolo attuale \\
    YearsSinceLastPromotion & int & [0, 15] & Anni dall'ultima promozione \\
    YearsWithCurrManager & int & [0, 17] & Anni con il manager attuale \\
\end{longtable}

\section{Spazio degli iperparametri e configurazione del modello}
\label{app:hyperparam}

La Tabella~\ref{tab:hyperparam_full} riporta lo spazio di ricerca completo esplorato dalla \texttt{RandomizedSearchCV} per il modello XGBoost. La configurazione della ricerca è la seguente:

\begin{itemize}
    \item Numero di iterazioni: 150
    \item Cross-validation: \texttt{StratifiedKFold} con 5 fold (\texttt{shuffle=True}, \texttt{random\_state=42})
    \item Metrica di scoring: F1-score sulla classe positiva
    \item Strategia di bilanciamento selezionata: \texttt{scale\_pos\_weight} (rapporto tra classe negativa e positiva sul train set)
    \item Seed di riproducibilità: \texttt{random\_state=42}
\end{itemize}

\begin{table}[htbp]
    \centering
    \caption{Spazio di ricerca completo degli iperparametri.}
    \label{tab:hyperparam_full}
    \begin{tabular}{lll}
        \toprule
        Iperparametro & Valori esplorati & Descrizione \\
        \midrule
        \texttt{n\_estimators}     & 100, 200, 300, 500, 700 & Numero di alberi \\
        \texttt{max\_depth}        & 3, 4, 5, 6, 7, 8 & Profondità massima per albero \\
        \texttt{learning\_rate}    & 0.01, 0.03, 0.05, 0.08, 0.1, 0.15 & Tasso di apprendimento \\
        \texttt{subsample}         & 0.6, 0.7, 0.8, 0.9, 1.0 & Frazione di campioni per albero \\
        \texttt{colsample\_bytree} & 0.5, 0.6, 0.7, 0.8, 0.9, 1.0 & Frazione di feature per albero \\
        \texttt{min\_child\_weight}& 1, 2, 3, 5, 7 & Peso minimo in un nodo foglia \\
        \texttt{gamma}             & 0, 0.05, 0.1, 0.2, 0.3 & Riduzione minima di loss per split \\
        \texttt{reg\_alpha}        & 0, 0.001, 0.01, 0.1, 1.0 & Regolarizzazione L1 \\
        \texttt{reg\_lambda}       & 0.5, 1.0, 1.5, 2.0, 3.0 & Regolarizzazione L2 \\
        \bottomrule
    \end{tabular}
\end{table}

Lo spazio combinatorio totale è pari a $5 \times 6 \times 6 \times 5 \times 6 \times 5 \times 5 \times 5 \times 5 = \num{3375000}$ configurazioni; la ricerca randomizzata ne ha campionate 150, corrispondenti allo 0{,}004\% dello spazio. Il modello finale, selezionato sulla base dell'F1 sul test set, è composto da 500 alberi.
